% Source: https://artofproblemsolving.com/community/c3h1646143_some_maths_problems_from_down_under
% Stable identity: aops-topic-1646143-q04
\textbf{If (3) is false,} then (1), (2) and (4) are all true. Two consecutive numbers $N-1$ and $N$ are both prime only for $N=3$ (one of them is even and greater than $2$ otherwise), and $3+2=5$ is not a square. So no $N$ works here.

\textbf{If (3) is true,} then $4<N<16$, and exactly two of (1), (2), (4) hold. Statement (4) holds in this range only for $N=7$ ($9-2$) and $N=14$ ($16-2$). If (4) fails, we would need $N$ and $N-1$ both prime, impossible for $N>3$. So (4) holds and exactly one of (1), (2) holds:
\begin{itemize}
\item $N=7$: $7$ is prime and $6$ is not. Exactly three true.
\item $N=14$: $14$ is not prime and $13$ is. Exactly three true.
\end{itemize}
The sum is $7+14=\boxed{21}$, as posted.
